% SLICECLUSTERMULT1 2026-10-18 03:15 Asia/Seoul
\documentclass[11pt]{article}
\usepackage[T1]{fontenc}
\usepackage{lmodern,amsmath,amssymb,amsthm,booktabs,hyperref}
\usepackage[margin=.9in]{geometry}
\newtheorem{theorem}{Theorem}
\newtheorem{lemma}[theorem]{Lemma}
\newtheorem{corollary}[theorem]{Corollary}
\newcommand{\EXACT}{\textbf{PROVED/EXACT}}
\newcommand{\KNOWN}{\textbf{KNOWN THEOREM}}
\newcommand{\COMP}{\textbf{COMPUTATIONAL}}
\newcommand{\CAND}{\textbf{CANDIDATE}}
\newcommand{\OPEN}{\textbf{OPEN}}
\newcommand{\RET}{\textbf{FALSIFIED/RETIRED}}
\title{SLICECLUSTERMULT1:\\Term-Count-Free Transfer for Arbitrary Bounded-Width DNFs via Spread-Cluster Multiplicity}
\author{P versus NP project}
\date{18 October 2026}
\begin{document}\maketitle

\noindent\textbf{Status.}
The canonical spread-cluster decomposition, exact-layer cluster multiplicity bound, conditional cluster-transfer theorem, and logarithmic-width transfer cone are \EXACT at the declared distributional models, pending independent review. The finite exhaustive/random audit is \COMP. CKR's robust-sunflower approximation and Rao's spread-matching method are \KNOWN. No circuit-to-DNF theorem with the required error profile is proved here; no unrestricted set-circuit lower bound or $P\ne NP$ proof is obtained.

\section{IDEAS PROPOSED BEFORE LITERATURE}
\paragraph{A: unique-core spread decomposition.}
Decompose an arbitrary bounded-width term family into spread residual families, then merge all leaves having the same core and residual width. Prediction: the number of clusters that can fire on one exact-$t$ input is bounded by a function of $t$ and width only, independent of the raw term count.

\paragraph{B: fully conditional cluster transfer.}
Condition simultaneously on a cluster core under the uniform exact slice and under the pseudorandom UNSAT distribution. Use two-sided relative moment matching to transfer residual terms and the exact-slice second-moment bound to remove the number of terms.

\paragraph{C: constant-width CKR repair.}
Ask whether the standard robust-sunflower approximation can approximate size-$N^K$ circuits using width $O_{K,B}(1)$ with inverse-polynomial error. Falsifier: the published gate-by-gate error is $\operatorname{size}(C)2^{-\Omega(c)}$, forcing $c=\Omega(\log N)$ for polynomial-size circuits.

\paragraph{D: logarithmic-width asymptotic cone.}
If constant width fails, retain width $W=\kappa\log q$ inside NAEGAUGEHIGHGIRTH1's moment window and characterize exactly how small the negative DNF acceptance must be to overcome cluster multiplicity.

\section{LITERATURE CHECK AND REVISIONS}
Cavalar--Kumar--Rossman's gate-by-gate approximator has total two-distribution error bounded by
\[
\operatorname{size}(C)\,2^{-\Omega(c)}.
\]
Thus their stock approximation requires $c=\Omega(\log N)$ against size $N^K$ circuits; Candidate C is not available from the published theorem without a new argument. Their closure/trimming framework remains valuable because it produces bounded-width minterm families.

Rao's spread-matching lower bound proves much stronger problem-specific monotone lower bounds, but its spread lemma necessarily has parameter dependence on both the object size and $\log(1/\varepsilon)$. The paper explicitly gives examples showing $\Omega(\log(1/\varepsilon))$ spread is necessary. Hence Rao's current theorem also does not supply a constant-parameter arbitrary-error approximation of our slice.

Kahn--Linial--Samorodnitsky show that low-order inclusion--exclusion data do not approximate arbitrary unions without structural hypotheses. Therefore a spread/cluster condition cannot be omitted.

The project already supplies the positive distribution needed below. NAEGAUGEHIGHGIRTH1 constructs, along an infinite sequence, an exact-density-$4/3$ UNSAT distribution $\mu_q^+$ with
\[
\Pr_{\mu_q^+}[A\subseteq X]=(1+o(1))\frac{(t)_{|A|}}{(M)_{|A|}}
\]
uniformly for every $|A|\le c\log q$ and every fixed $c<c_*$, where
\[
c_*=\frac{2}{3\log11+2\log(13/4)}=0.20940225\ldots.
\]
This run treats that theorem as the moment input and develops the missing arbitrary-DNF transfer.

\section{CONSTRUCTION OR PROOF ATTEMPT}
Let $\nu$ be the uniform exact-$t$ distribution on $\binom{[N]}t$. A monotone term family $\mathcal F$ of width at most $W$ represents
\[
D_{\mathcal F}(x)=\bigvee_{T\in\mathcal F}x_T,\qquad x_T=\prod_{i\in T}x_i.
\]

\subsection{Canonical spread-cluster decomposition}
A $w$-uniform family $\mathcal G$ is $\rho$-spread if
\[
d_{\mathcal G}(A):=|\{R\in\mathcal G:A\subseteq R\}|
\le |\mathcal G|\rho^{-|A|}
\]
for every $A$.

A cluster is a triple $(C,w,\mathcal G)$ representing the terms
\[
\{C\cup R:R\in\mathcal G\},
\]
where each $R$ is disjoint from $C$, has size $w$, and $\mathcal G$ is $\rho$-spread.

\begin{theorem}[Unique-core spread decomposition]
For every finite term family $\mathcal F$ of width at most $W$ and every $\rho\ge1$, there is a partition of $\mathcal F$ into $\rho$-spread clusters with at most one cluster for each pair $(C,w)$ satisfying
\[
|C|+w\le W.
\]
\end{theorem}
\begin{proof}
First partition $\mathcal F$ by term size. For a uniform family $\mathcal H$, if it is $\rho$-spread, output it. Otherwise choose a nonempty witness $A$ with
\[
d_{\mathcal H}(A)>|\mathcal H|\rho^{-|A|}.
\]
Partition $\mathcal H$ into terms containing $A$ and terms not containing $A$. Strip $A$ from the former and add it to the core; recurse on both branches. Core size strictly increases on every containing branch, while every partition strictly reduces unresolved term mass, so the recursion terminates.

The leaves partition the original terms. Finally merge all leaves having the same core $C$ and residual width $w$. The union of $\rho$-spread $w$-uniform families is $\rho$-spread because both codegrees and family sizes add. Thus there is at most one cluster per $(C,w)$.
\end{proof}

\subsection{One-input cluster multiplicity}
For an exact-$t$ input $X$, a cluster can fire only if its core is contained in $X$.

\begin{theorem}[Cluster multiplicity]
For every exact-$t$ input $X$ and every width-$W$ term family,
\[
\#\{\text{clusters firing on }X\}
\le
\mathfrak M_W(t)
:=
\sum_{c=0}^W (W-c+1)\binom tc.
\]
Consequently, if $E_j$ are the cluster events under any distribution supported on the exact layer,
\[
\sum_j\Pr[E_j]\le \mathfrak M_W(t)\Pr\!\left[\bigcup_jE_j\right].
\]
\end{theorem}
\begin{proof}
For core size $c$, at most $\binom tc$ cores can be subsets of $X$, and for each core there are at most $W-c+1$ possible residual widths. This proves the pointwise multiplicity bound. Taking expectations of
\[
\sum_j1_{E_j}(X)\le \mathfrak M_W(t)1_{\cup_jE_j}(X)
\]
proves the probability inequality.
\end{proof}

This removes the global cluster-count parameter. An arbitrary number of mass-produced terms may exist; only clusters simultaneously compatible with one exact-layer input are charged.

\subsection{Conditional spread transfer}
Let $\mu$ be another distribution on the exact layer and assume two-sided relative moment matching through order $W$:
\[
(1-\delta)\nu[A]\le\mu[A]\le(1+\delta)\nu[A],
\qquad |A|\le W,
\]
where $0\le\delta<1$.

Fix one cluster $(C,w,\mathcal G)$ and put
\[
c=|C|,\quad
N_c=N-c,\quad
t_c=t-c,\quad
p_c=t_c/N_c.
\]
Conditioned on $C\subseteq X$, $\nu$ is exactly the uniform $t_c$-slice on the remaining $N_c$ coordinates. Let
\[
\eta_C=\Pr_{\nu}[D_{\mathcal G}=1\mid C],
\quad
c_0=\exp\!\left(-\frac{w(w-1)}{t_c}\right),
\quad
a=\frac{w}{\rho p_c},
\quad
r_\delta=\frac{1+\delta}{1-\delta}.
\]

\begin{theorem}[Conditional cluster transfer]
If $w\le t_c/2$, then
\[
\Pr_\mu[C\wedge D_{\mathcal G}]
\le
K(C,w)\,
\Pr_\nu[C\wedge D_{\mathcal G}],
\]
where
\[
K(C,w)
=
\frac{2(1+\delta)}{c_0^2}
\max\{r_\delta,e^a\}.
\]
\end{theorem}
\begin{proof}
For every residual $A$ with $|C|+|A|\le W$,
\[
\Pr_\mu[A\mid C]
=
\frac{\Pr_\mu[A\cup C]}{\Pr_\mu[C]}
\le
r_\delta\,\Pr_\nu[A\mid C].
\]
The exact-slice spread-DNF theorem applied conditionally gives
\[
\eta_C\ge\frac{c_0^2\lambda_C}{1+e^a\lambda_C},
\qquad
\lambda_C=|\mathcal G|p_c^w.
\]
If $\eta_C\le c_0^2/(2e^a)$, rearranging yields
\[
\lambda_C\le 2\eta_C/c_0^2.
\]
A union bound and conditional relative moments give
\[
\Pr_\mu[D_{\mathcal G}\mid C]
\le r_\delta\lambda_C
\le 2r_\delta\eta_C/c_0^2.
\]
If $\eta_C>c_0^2/(2e^a)$, the trivial upper bound $1$ is at most
\[
2e^a\eta_C/c_0^2.
\]
Multiplying by $\Pr_\mu[C]\le(1+\delta)\Pr_\nu[C]$ proves the claim.
\end{proof}

\subsection{Arbitrary bounded-width DNF transfer}
Combine the preceding two theorems.

\begin{corollary}[Term-count-free bounded-width transfer]
Assume the moment condition through order $W$, and choose $\rho$ so that
\[
\inf_{0\le c\le W}\rho\,\frac{t-c}{N-c}\ge\theta>0.
\]
Then every monotone DNF of width at most $W$ satisfies
\[
\boxed{
\Pr_\mu[D_{\mathcal F}=1]
\le
K_W(\delta,\theta,t)\,
\mathfrak M_W(t)\,
\Pr_\nu[D_{\mathcal F}=1],
}
\]
where
\[
K_W
=
2(1+\delta)
\exp\!\left(\frac{2W(W-1)}{t-W}\right)
\max\!\left\{\frac{1+\delta}{1-\delta},e^{W/\theta}\right\}.
\]
The bound is independent of $|\mathcal F|$.
\end{corollary}

For fixed $W$ and $\delta=o(1)$, this costs only $O_W(t^W)$. For logarithmic width it yields an explicit asymptotic cone.

\subsection{Logarithmic-width transfer cone}
Let
\[
t=\Theta(q),\qquad W=\kappa\log q,\qquad 0<\kappa<c_*.
\]
Choose $\rho=\theta/p$ with fixed $\theta>0$ and $p=t/N$. Since $W=o(t)$,
\[
\log K_W=O(W)=O(\log q).
\]
Moreover
\[
\mathfrak M_W(t)
\le
(W+1)^2\left(\frac{et}{W}\right)^W,
\]
and hence
\[
\boxed{
\log\mathfrak M_W(t)
=
\kappa(\log q)^2
-\kappa\log q\log\log q
+O_\kappa(\log q).
}
\]

\begin{corollary}[NAE logarithmic-width invisibility cone]
Let $\mu_q^+$ be the NAEGAUGEHIGHGIRTH1 UNSAT distribution and $\nu_q$ the uniform exact-density-$4/3$ layer. Fix
\[
0<\kappa<c_*.
\]
For every width-$W_q\le\kappa\log q$ monotone DNF,
\[
\boxed{
\Pr_{\mu_q^+}[D=1]
\le
\exp\!\left((\kappa+o(1))(\log q)^2\right)
\Pr_{\nu_q}[D=1].
}
\]
Therefore, if for some fixed $\varepsilon>0$
\[
\Pr_{\nu_q}[D=1]
\le
\exp\!\left(-(\kappa+\varepsilon)(\log q)^2\right),
\]
then
\[
\Pr_{\mu_q^+}[D=1]=o(1).
\]
\end{corollary}

This is the main new conclusion. It applies to \emph{arbitrary} monotone DNFs inside the logarithmic moment window, regardless of their number of terms.

\subsection{What is still needed for circuits}
The stock CKR and Rao frameworks do not yet provide the required final approximation. Their error scales force a logarithmic approximation parameter for polynomial-size circuits, which is compatible with the moment window, but the present DNF transfer then requires negative acceptance on the scale
\[
\exp(-\Theta((\log q)^2)),
\]
not merely $q^{-O(1)}$.

A sufficient final circuit theorem can therefore be stated cleanly.

\begin{definition}[Required approximation profile]
For each fixed polynomial circuit exponent $K$, find constants
\[
0<\kappa_K<c_*,\qquad \varepsilon_K>0
\]
such that every size-$M^K$ monotone circuit computing the NAE slice has an approximating monotone DNF $D$ of width at most $\kappa_K\log q$ satisfying
\[
\Pr_{\mu_q^+}[C=1,D=0]=o(1)
\]
and
\[
\Pr_{\nu_q}[D=1]\le
\exp\!\left(-(\kappa_K+\varepsilon_K)(\log q)^2\right).
\]
\end{definition}

If this profile holds, the logarithmic-width invisibility corollary makes $\Pr_{\mu_q^+}[D=1]=o(1)$, contradicting the fact that $C=1$ on every positive input.

This avoids selecting one universal numerical parameter tuple. It is an open asymptotic region indexed by the hypothetical circuit exponent $K$.

\section{TESTS AND OUTCOME}
The dependency-free audit used seed 20261018 and tested random term families in three exact-slice regimes:
\[
(N,t,W,\rho)=(8,3,2,2),\ (9,4,3,2),\ (10,4,3,2.5).
\]
For each regime it generated 40 random families, recursively decomposed and merged them, and exhaustively enumerated all exact-layer inputs.

It checked:
\begin{itemize}
\item exact reconstruction of every original term family from the clusters;
\item uniqueness of every $(C,w)$ key;
\item spreadness of every residual family;
\item the pointwise cluster-multiplicity bound;
\item $\sum_j\Pr_\nu[E_j]\le\mathfrak M_W(t)\Pr_\nu[D]$;
\item the conditional exact-slice spread second-moment lower bound whenever its hypotheses applied.
\end{itemize}

There were 17,385 declared assertions across 120 test families, zero failures, at most 19 produced clusters, and at most 7 simultaneously firing clusters. Script SHA-256:
\[
\texttt{93c6dcce0762ddcf52df379be8b74c2595dab7f45c1975be1c7657d5dacdfd6c}.
\]
Raw-result SHA-256:
\[
\texttt{a1ace045db8c3c7d46f24a64693e0fb3bbce62de448b01d1c543b7422aa36ac5}.
\]

\paragraph{Outcome.}
Candidate A is \EXACT and stronger than originally proposed: raw cluster count disappears in favor of exact-layer firing multiplicity. Candidate B is \EXACT. Candidate C is \RET in the stock CKR/Rao forms because polynomial circuit size forces a logarithmically growing approximation parameter. Candidate D survives in a precise form: logarithmic width is compatible with the existing NAE moment window, but it demands quasi-polynomially small negative acceptance.

\paragraph{Next target.}
The highest-value next step is to remove the remaining $\mathfrak M_W(t)$ loss by exploiting \emph{weighted core rarity} rather than worst-case firing multiplicity. Search for a canonical dense-core decomposition satisfying
\[
\sum_{\text{clusters}}\Pr_\nu[C]
\le q^{O(W)}\quad\text{or ideally }\exp(O(W\log W)),
\]
with the residual transfer charged relative to the conditional cluster acceptance rather than one per possible core. A bound $\exp(O(W\log W))$ would turn $W=\Theta(\log q)$ into only $\exp(O(\log q\log\log q))$ loss and would reduce the needed negative error from $\exp(-\Theta((\log q)^2))$ to quasi-polynomially milder scale.

In parallel, construct a SAT-supported companion distribution to $\mu_q^+$ with the same gauge/high-girth local moments. Such a pair would remove the residual contribution of rare UNSAT formulas from the uniform negative layer and is the cleanest interface for CKR-style two-distribution approximation.

Root \texttt{main.tex} remains truncated and was not modified.

\begin{thebibliography}{4}
\bibitem{CKR22} B. P. Cavalar, M. Kumar, B. Rossman, \emph{Monotone Circuit Lower Bounds from Robust Sunflowers}, Algorithmica 84, 2022.
\bibitem{Rao26} A. Rao, \emph{Monotone Circuit Lower Bounds from Spread Matchings}, ECCC TR26-129, 2026.
\bibitem{KLS96} J. Kahn, N. Linial, A. Samorodnitsky, \emph{Inclusion-exclusion: Exact and approximate}, Combinatorica 16, 1996.
\bibitem{NAE} Project report NAEGAUGEHIGHGIRTH1, 17 October 2026.
\end{thebibliography}
\end{document}
% FINAL_VERIFICATION_SENTINEL::MAIN::SLICECLUSTERMULT1::2026-10-18_0315::P_VS_NP_OPEN
